\documentclass[10pt]{article}
\usepackage{template/template}
\NewsSetKeys{date={2026年11月6日},vol={024},title={每日有趣新闻},author={编辑部},note={本期短刊：资料截至北京时间11月6日上午。数学新稿已检索，但尚未完成原始预印本核验。}}
\begin{document}
\begin{documents}
\newsitem{空间站最后一组升级太阳翼完成安装}{航天｜11月5日｜NASA、Spaceflight Now}{NASA宇航员Anil Menon和加拿大宇航员Joshua Kutryk于11月5日出舱，为国际空间站安装最后一组可卷展太阳翼iROSA。这项持续数年的供电升级接近完成。宇航员利用机械臂转运设备，固定支架，连接电缆并展开太阳翼。任务期间出现供电连接问题，经过地面与舱外人员协同排查后恢复供电。可卷展太阳翼的关键在于发射时收纳紧凑、在轨展开获得较大的光照面积；安装还必须适应轨道昼夜周期和舱外活动安全限制。NASA发布了任务安排，Spaceflight Now提供现场进程记录。安装完成并不意味着全部在轨性能测试已经结束。\cite{nasa,space}}
\newsitem{以色列组阁：推荐人数不是执政联盟}{国际政治｜11月5日｜路透社}{以色列10月27日选举未产生明确执政多数。11月5日，120名议员中67人推荐加迪·艾森科特尝试组阁，46人推荐内塔尼亚胡。总统赫尔佐格表示周日决定提名人选，同时指出党派间存在相互否决关系。议员支持某人获得组阁机会，不代表愿意加入其政府。能否形成多数执政联盟，仍取决于之后的谈判与协议。因此不能将67票推荐误写成已经组建67席政府。\cite{israel}}
\newsitem{农产品采购：谈判承诺与进口统计不可混同}{中美贸易｜11月5日｜路透社}{中国方面表示，扩大采购美国农产品取决于互惠关税削减措施落实。此前美国政府提及约170亿美元年度农产品采购安排，但具体执行时间和条件仍有分歧。报道指出今年1至9月非大豆农产品进口额约72亿美元。承诺采购金额并不等于已经完成的海关进口金额，实际贸易还受关税、运输与市场价格影响。后续需要同时跟踪正式协议与实际到货数据。\cite{farm}}
\newsitem{GDP重估：埃塞俄比亚经济规模为何增加近六成？}{经济与统计｜11月5日｜路透社}{埃塞俄比亚更新国民账户基期后，将2024／25年度GDP规模上修59.2\%，至约2014.2亿美元。这是统计基期和经济活动覆盖范围调整产生的重新估计，并非一夜之间创造了六成的新产出。重估会改变债务与GDP比值的分母，却不会自动减少外币债务本金，因此不能据此推断偿债能力同比例提升。\cite{ethiopia}}
\newsitem{数学检索说明}{数学｜11月5—6日}{已检索数论、代数、几何、组合、逻辑、概率及Lean/mathlib新研究；尚未取得足以独立核对作者、版本与定理陈述的当日原始预印本。本期不刊登未经核查的数学突破。}
\newsitem{出版说明}{编辑部}{因材料与版面制作时间限制，本期为短刊，不冒充符合常规深稿与现场配图要求的完整报纸。}
\bibliographystyle{plain}\bibliography{refs}
\end{documents}
\end{document}